\section{近期研究（2026 年 8--11 月）}
\paragraph{G-Skin：以生成姿态图像为高斯骨架绑定提供监督}
2026.08~\cite{2608_01726}。G-Skin 接收已有静态高斯和给定骨架，将二维生成模型产生的多姿态图像用作蒙皮监督，在测试对象上优化高斯随骨骼运动的权重。几何正则约束相邻位置的变形，减弱生成图像不准确时的绑定漂移。

作者用 50 个类人和 35 个非人形的生成资产评价。关联部位已经粘连或道具与身体拓扑融合时，局部连续性可能促成错误绑定。此流程仍需逐对象优化，不能称为一步预测骨架或蒙皮权重。

\paragraph{PD-GS：把离散音素和连续语音融合到高斯头像动画}
2026.08~\cite{2608_05218}。PD-GS 先建立规范姿态的高斯头部，再用连续音频和对齐的音素共同生成形变。融合模块在嘴唇闭合等发音关键区段更重视离散语言线索，连续音频仍提供上下文。

作者在 HDTF 的主要评价协议中使用两个常见对象，分别比较唇部关键点、同步、渲染质量和模型侧延迟。这属于个体特定头像实验，少样本人脸推广及舌等其他发音器官仍是后续问题。

\paragraph{4DAnyone：用骨架条件与跨视图合成支持单目人体四维重建}
2026.08~\cite{2608_20335}。4DAnyone 从单目视频恢复人体骨架，用骨架作为稀疏几何条件生成多视角视频。参考上下文压缩减少视角增加时的外观条件开销，目标视图的分组在去噪中动态关联，以减小各组之间的结构漂移。生成结果随后用于高斯人体重建。

测试覆盖 DNA-Rendering 和 DyMVHumans，并检查多视角合成与下游四维高斯结果。论文的失败例显示，远离骨架的宽松衣物可能跨视图不一致，错误骨架也可能被一致地传播到全部新视角。

\paragraph{ManifoldSplat：用自然语言改变绑定到高斯头像的局部几何}
2026.10~\cite{2610_03599}。ManifoldSplat 从单目视频建立可动画头像，根据文本生成局部 FLAME 网格形状增量，并传递给绑定于网格的高斯。几何编辑和渲染保持不同角色，表面一致性约束用于抑制形状改变后原语与网格脱节。

作者在 NeRSemble 子集保留正面相机，并用提示一致性、身份和渲染指标评价多种局部编辑。大形变会放大 FLAME 与原高斯的残余不匹配，实际编辑空间也受预定义形变原型限制。

\paragraph{HumanSkill-4D：用持续语义区域和关节轨迹回答头像查询}
2026.10~\cite{2610_24452}。HumanSkill-4D 把人体部位与衣物标签关联到已有码本中的持久高斯区域，让固定多模态模型能按语义取出原语。涉及动作时，模型把自然语言条件组合为关节轨迹上的运动和时间算子，返回匹配区域及时间段。流程不对单个头像另训练。

作者在扩展已有动作和人体数据的三个指代任务上评价。冻结解析器的类别和动作模板有明确边界，例如未登记的小饰品或动作无法仅凭自由文本可靠查询。

\paragraph{EGGHead：在固定原语预算内重建单图头部高斯}
2026.10~\cite{2610_26898}。EGGHead 学习按头部语义区域分配有限高斯，并借助 Voronoi 划分推断稀疏原语的局部形状。轻量对称模块补全遮挡区域，推断无需为每个新对象再执行逐个优化。

论文在不重合的 FFHQ 和 CelebA-HQ 图像上评价。模型未将纹理拆成反照率与照明，不能用于重打光；不均匀光线会影响对称补全，完全正脸时后脑仍可能不完整。

\paragraph{World2HOSI：让长流程人体与物体交互共享可更新状态}
2026.10~\cite{2610_27725}。World2HOSI 将静态高斯外观、可交互几何和持久物体身份组织成统一场景状态。导航、人与物体交互和人体与场景互动由不同专家合成，每段动作完成后的物体变化再写入状态，为下一段确定位置与目标。

论文在十个真实重建或合成场景的二十条日常指令上度量任务完成和几何质量，并展示有跨段依赖的动作。可选物理修正以专家给出的参考动作为前提，难以纠正本身不可行的大偏差计划。

\paragraph{复杂度先验引导通用人体高斯的分配与局部监督}
2026.10~\cite{2610_35152}。该方法从衣着人体网格的局部几何与纹理变化估计顶点复杂度，据此把更多高斯分配给褶皱或细节区，并在训练时更经常对这些区域抽取局部图像监督。整体图像监督仍保留，以防只优化高频区域。

作者用 THuman2.0、2K2K 和 X-Humans 的已标定渲染视图比较。这是静态多视图人体重建方法，逐帧应用到视频并不会自带时间一致性，也不能与单张无标定头像重建共享同一输入口径。
